\begin{proof}
    Let $\alpha \in \mathbb{R}$.
    First, it must be shown that $U \odot e^{i\alpha\pi} = U^\alpha$ for all $U \in \Mor( \CatUnitary )$ if and only if $\alpha \in (-1, 1]$.
    \begin{itemize}
    \item Assume that $\alpha \in ( -1, 1 ]$.
          Then $i\alpha\pi \in i( -\pi, \pi ]$.
          Then $\Log( e^{i\alpha\pi} ) = i\alpha\pi$ since $\Log( - )$ is the principle branch of $\log( - )$.
          Let $U \in \Mor( \CatUnitary )$.
          Then the following equation holds.
          \begin{equation*}
              U \odot e^{i\alpha\pi}
              =
              \exp\left( \Log( U ) \otimes ( i\alpha\pi ) / (i\pi) \right)
              =
              \exp( \alpha \Log( U ) )
              =
              U^\alpha
          \end{equation*}
          Since $U$ was arbitrary, then $U \odot e^{i\alpha\pi} = U^\alpha$ for all $U \in \Mor( \CatUnitary )$.
    \item Assume that $\alpha \not \in ( -1, 1 ]$.
          Then there exists a unique $k \in \mathbb{Z} \setminus \{ 0 \}$ such that $\alpha - 2k \in ( -1, 1 ]$.
          Then $i\alpha\pi - i2k\pi \in i( -\pi, \pi ]$, which implies that $\Log( e^{i\alpha\pi} ) = i\alpha\pi - i2k\pi$.
          Since $k \ne 0$, then let $U = Z^{1/|4k|}$, so that $\Log( U ) = i \frac{\pi}{|4k|} \ket{1} \bra{1}$.
          Then the following equations hold.
          \begin{align*}
              U^\alpha
              &=
              \exp\left( \alpha \left( i\frac{\pi}{|4k|} \ket{1} \bra{1} \right) \right)
              =
              \exp\left( i \frac{\alpha\pi}{|4k|} \ket{1} \bra{1} \right)
              \\
              U \odot e^{i\alpha\pi}
              &=
              \exp\left( \left( i\frac{\pi}{|4k|} \ket{1} \bra{1} \right) \otimes ( i\alpha\pi - i2k\pi ) / (i\pi) \right)
              =
              \exp\left( i \frac{(\alpha-2k)\pi}{|4k|} \ket{1} \bra{1} \right)
          \end{align*}
          Since $-1 < \alpha - 2k \le 1$, then $2k - 1 < \alpha < 2k + 1$.
          The value of $\alpha / |4k|$ then depends on the value of $k$.
          If $k > 0$, then the following equations hold.
          \begin{align*}
              (2k - 1)/|4k| &= 2k/|4k| - 1/|4k| = 1/2 - 1/|4k| \ge 1/2 - 1/4 > -1
              \\
              (2k + 1)/|4k| &= 2k/|4k| + 1/|4k| = 1/2 + 1/|4k| \le 1/2 + 1/4 < 1
          \end{align*}
          This means that $\alpha / |4k| \in ( -1, 1 )$.
          If $k < 0$, then the following equations hold.
          \begin{align*}
              (2k - 1)/|4k| &= 2k/|4k| - 1/|4k| = -1/2 - 1/|4k| \ge -1/2 - 1/4 > -1
              \\
              (2k + 1)/|4k| &= 2k/|4k| + 1/|4k| = -1/2 + 1/|4k| \le -1/2 + 1/4 < 1
          \end{align*}
          This means that $\alpha / |4k| \in ( -1, 1 )$.
          Then in either case, $\alpha / |4k| \in ( -1, 1 )$.
          It follows that $\Log( U^\alpha ) = i \frac{\alpha\pi}{|4k|} \ket{1} \bra{1}$.
          Since $i( \alpha - 2k )\pi \in i( -\pi, \pi ]$, then $i \frac{(\alpha - 2k)\pi}{|4k|} \in i( -\pi, \pi ]$.
          It follows that $\Log( U \odot e^{i\alpha\pi} ) = i \frac{(\alpha - 2k)\pi}{|4k|} \ket{1} \bra{1}$.
          Clearly $\Log( U^\alpha ) \ne \Log( U \odot e^{i\alpha\pi} )$.
          Then $U^\alpha \ne U \odot e^{i\alpha\pi}$.
          Then there exists a $U \in \Mor( \CatUnitary )$ such that $U \odot e^{i\alpha\pi} \ne U^\alpha$.
    \end{itemize}
    It remains to be shown that if $U \in \Mor( \CatUnitary )$ is Hermitian, then $U \odot e^{i\alpha\pi} = U^\alpha$.
    Assume that $U \in \Mor( \CatUnitary )$ with $U$ Hermitian.
    Let $\{ \beta_j \}_{j=1}^n$ denote the eigenvalues of $\Log( U ) / (i\pi)$ with orthonormal eigenbasis $\{ b_j \}_{j=1}^n$.
    Then the following equation holds where $z = e^{i\alpha\pi}$.
    \begin{equation*}
        U \odot z
        =
        \exp\left( \Log( U ) \otimes \Log( z ) / (i\pi) \right)
        =
        \exp\left( \sum_{j=1}^n \Log( z ) \beta_j \ket{b_j} \bra{b_j} \right)
        =
        \sum_{j=1}^n e^{\Log( z ) \beta_j} \ket{b_j} \bra{b_j}
    \end{equation*}
    Let $j \in \{ 1, 2, \ldots, n \}$.
    Then $\exp( i\beta_j\pi )$ is an eigenvalue of $U$.
    Since $U$ is Hermitian, then $e^{i\beta_j\pi}$ is either $-1$ or $1$.
    Then $\beta_j\pi$ is a multiple of $\pi$.
    Since $\beta_j\pi \in ( -\pi, \pi ]$, then $\beta_j$ is either $0$ or $1$.
    If $\beta_j = 0$, then the following equation holds.
    \begin{equation*}
        \exp( \Log(z)\beta_j )
        =
        \exp( 0 \Log(z) )
        =
        \exp( 0 )
        =
        \exp( 0 (i\alpha\beta_j) )
        =
        \exp( i\alpha\beta_j\pi )
    \end{equation*}
    If $\beta_j = 1$, then the following equation holds.
    \begin{equation*}
        \exp( \Log(z)\beta_j )
        =
        \exp( 1 \Log(z) )
        =
        \exp( \Log(z) )
        =
        \exp( i\alpha\pi )
        =
        \exp( 1 (i\alpha\pi) )
        =
        \exp( i\alpha\beta_j\pi )
    \end{equation*}
    In either case, $e^{\Log(z)\beta_j} = e^{i\alpha\beta_j\pi}$.
    Since $j$ was arbitrary, then the following equation holds.
    \begin{align*}
        \sum_{j=1}^n e^{\Log(z)\beta_j} \ket{b_j} \bra{b_j}
        =
        \sum_{j=1}^n e^{i\alpha\beta_j\pi} \ket{b_j} \bra{b_j}
        =
        \exp\left( \sum_{j=1}^n i\alpha\beta_j\pi \ket{b_j} \bra{b_j} \right)
        =
        \exp( \alpha \Log( U ) )
        =
        U^\alpha
    \end{align*}
    In conclusion, $U \odot e^{i\alpha\pi} = U^\alpha$.
    Since $U$ was arbitrary, then $U \odot e^{i\alpha\pi} = U^\alpha$ for each Hermitian $U \in \Mor( \CatUnitary )$.
    Since $\alpha$ was arbitrary, then this completes the proof.
\end{proof}